% META: {"id":"1SPE-SUITES-CO-052","chapitre":"1SPE-SUITES","type_objet":"corrige","exercice_id":"1SPE-SUITES-EX-052","status":"generated"}
% BEGIN-VERIFY
% from math import isclose
% assert isclose((250 - 200) / 200 * 100, 25.0, abs_tol=1e-12)
% c = 1.25
% cr = 1 / c
% assert isclose(cr, 0.8, abs_tol=1e-12)
% assert isclose((cr - 1) * 100, -20.0, abs_tol=1e-12)
% # Une baisse de 25 % apres une hausse de 25 % ne rend pas le prix initial.
% assert isclose(250 * 0.75, 187.5, abs_tol=1e-12)
% assert 187.5 != 200
% # n hausses suivies de n baisses reciproques : le produit vaut 1.
% for n in range(1, 6):
%     assert isclose(200 * c**n * cr**n, 200, abs_tol=1e-9)
% END-VERIFY
\begin{corrige}{1SPE-SUITES-EX-052}
\begin{enumerate}
\item Le taux d'évolution vaut
  \[
    \frac{250 - 200}{200} = 0{,}25, \qquad\text{soit } +25\,\%.
  \]

\item Passer de $250$ à $200$ revient à multiplier par
  $\dfrac{200}{250} = 0{,}8$. Le coefficient réciproque est donc l'inverse du
  premier : $\dfrac{1}{1{,}25} = 0{,}8$. Le taux correspondant est
  $0{,}8 - 1 = -0{,}2$, soit $-20\,\%$.

\item Une baisse de $25\,\%$ appliquée à $250$ euros donne
  $250 \times 0{,}75 = 187{,}50$ euros, et non $200$. La baisse porte sur une
  valeur plus grande que la valeur initiale : le taux réciproque d'une hausse
  de $25\,\%$ est $-20\,\%$, pas $-25\,\%$.

\item Après $n$ hausses, le prix vaut $u_n = 200 \times 1{,}25^{\,n}$. Chaque
  baisse réciproque multiplie par $0{,}8 = \dfrac{1}{1{,}25}$, donc après $n$
  baisses :
  \[
    200 \times 1{,}25^{\,n} \times 0{,}8^{\,n}
    = 200 \times (1{,}25 \times 0{,}8)^{\,n}
    = 200 \times 1^{\,n} = 200.
  \]
  Le prix revient exactement à sa valeur initiale : c'est ce que signifie
  « réciproque ».
\end{enumerate}
\end{corrige}
